% ------------------------------------------------------------------------------
% 11.1. IMPLEMENTACIÓN DEL SISTEMA
% ------------------------------------------------------------------------------
\subsection{11.1. Implementación del Sistema}\label{sec:cap11_implementacion}

El Sistema ComedorUAGRM busca lograr los siguientes resultados:

\begin{enumerate}
    \item \textbf{Digitalización de la convocatoria y la postulación:}
    \begin{itemize}
        \item Postulación 100~\% en línea, desde el panel web o la aplicación móvil, con digitalización de los documentos mediante la cámara del dispositivo.
        \item Validación de campos y documentos obligatorios, de la elegibilidad y de los duplicados antes de enviar la solicitud.
        \item Consulta del estado de la postulación por parte del postulante, con avisos por cada cambio de estado.
    \end{itemize}

    \item \textbf{Objetividad en la preselección y la adjudicación:}
    \begin{itemize}
        \item Cálculo del puntaje de vulnerabilidad con inteligencia artificial para todas las solicitudes recibidas, y no solo para las que se alcanzan a revisar.
        \item Lista priorizada que permite a la Dirección Universitaria de Bienestar Social y Salud (DUBSS) auditar únicamente a los finalistas mejor puntuados.
        \item Adjudicación respaldada por un acta y por el registro del puntaje con su fecha de cálculo, lo que da trazabilidad a cada decisión.
    \end{itemize}

    \item \textbf{Control de becados y de la asistencia al comedor:}
    \begin{itemize}
        \item Registro digital diario de la asistencia por parte del personal del comedor.
        \item Generación de planillas de asistencia imprimibles y cotejo del registro digital con el respaldo físico.
        \item Gestión del padrón de becados (baja, suspensión y reemplazo) y consulta del historial de asistencia por el propio becado.
    \end{itemize}

    \item \textbf{Información para la toma de decisiones:}
    \begin{itemize}
        \item Dashboard de indicadores de Business Intelligence para los directivos de la DUBSS.
        \item Bitácora de acciones que permite consultar quién realizó cada operación relevante y cuándo.
    \end{itemize}

    \item \textbf{Fortalecimiento de la seguridad:}
    \begin{itemize}
        \item Control de accesos mediante roles y permisos específicos para el postulante, el administrador de la DUBSS, el directivo y el personal del comedor.
        \item Autenticación con tokens JWT y almacenamiento de las contraseñas como hash.
        \item Protección de los datos personales y de los documentos de los postulantes mediante el almacenamiento controlado de archivos.
    \end{itemize}
\end{enumerate}
